\section{The Interoperability Problem}

RCOUNT is only useful if it can meet election offices where their public records
already are. Some offices publish machine-readable result exports. Others
publish statements of votes as spreadsheets, PDFs, precinct reports, county
canvass files, or vendor-specific artifacts. The EAC describes canvass and
certification as the process by which unofficial results become official public
results \citep{eac-canvass-certification}. RCOUNT's role is narrower: it
normalizes declared public evidence into a package that can be mechanically
rechecked.

The adapter layer is therefore not a cosmetic import step. It is where source
meaning becomes explicit. A field named ``total'' might mean votes for one
selection, ballots cast in a contest, ballots accepted in a batch, or a
jurisdiction-wide summary. A status label might mean election-night unofficial,
canvassed, recounted, amended, or certified. An adapter must not guess silently.
It must map each source field to a typed RCOUNT object or leave the source claim
outside the verifier.
